\section*{Appendix}

\section{Prompt Details}
\subsection{For Concatenation}
\label{subsec:prompt_concat}
Table \ref{table:appendix_concatenation_prompt} depicts the detailed prompts of concatenating multiple utterances using ChatGPT. 
When two or three sentences are given as input, the prompt works to combine them into a single sentence while maintaining the essence of each sentence, mirroring the exemplary good answers. 
Despite numerous constraints, the generative approach often struggles to achieve concatenation, as shown in Table \ref{table:generative}.
We experimented with a variety of prompts, including adding or excluding components, in an effort to optimize the concatenation process. 
The term [Bad Answer] indicates an answer is not recommended, not necessarily incorrect.
However, the variations in prompts yielded negligible differences in the quality of the results.

\begin{table*}
    \small
    \centering
    \renewcommand{\arraystretch}{1.5}
    \setlength{\tabcolsep}{0.5em}
    
    \begin{tabular}{>{\raggedright\arraybackslash}p{15cm}}
        \toprule
        You are a native English speaker. \\
        
        \textbf{[Task Definition]} Combine 2 or 3 utterances as one single utterance.\\
        
        \textbf{[Goal]} The focus is on creating a single utterance that captures the essence of both ideas without unnecessary redundancy.\\
        
        \textbf{[Instructions]}
        \begin{itemize}[leftmargin=10mm, nosep]
            \item Avoid adding just punctuation.
            \item Don't paraphrase.
            \item Don't compromise the meaning of each utterance.
            % \item Don't capitalize all characters.
            \item Don't replace numbers with radix.
            \item Maintain the intent of each utterance.
            \item Don't forget that if a utterance starts with a verb, it's a statement.
            \item Do NOT use conjunctions like `and'.
            % \item Don't print `[Good Answer]'.
            \item Don't print intent directly. 
        \end{itemize} \\
        \arrayrulecolor{lightgray}\midrule\arrayrulecolor{black} 
        
        \textcolor{commentcolor}{\# $N$ iterations}\\
        \textbf{[Example 1]} 
        % \textbf{[Example $N$]} 
        \begin{quote}
            
        play my 88 keys playlist (\texttt{PlayMusic}) + add another song to my 88 keys playlist (\texttt{AddToPlaylist}) 
        
        {[Good Answer]} while playing my 88 keys playlist, add another song to it. 
        
        {[Bad Answer]} Play my 88 keys playlist and also add another song to my 88 keys playlist. 
        
        \end{quote}
        ... \\
        
        \arrayrulecolor{lightgray}\midrule\arrayrulecolor{black}
        
        \textbf{[Query]} Combine the following utterances naturally.
        
        Inside the parentheses is the intent of each utterance: $\mathrm{utt_0}$ ({$\mathrm{intent_0}$}) + $\mathrm{utt_1}$ ({$\mathrm{intent_1}$}) \\
        
        \bottomrule
    \end{tabular}
    
    \caption{Specification of the prompt used for the Generative Concatenation Approach with $N=3$.}
    \label{table:appendix_concatenation_prompt}
\end{table*}


\subsection{For Evaluating Multi-intent Detection}
\label{subsec:prompt_evaluate}
Table \ref{table:appendix_chatgpt_prompt} is the prompt for evaluating MID using ChatGPT in our study. Each dataset offers its distinct set of labels and examples to detect multiple intents of the given query.
In this study, we expanded the dataset up to 150 single intents into multi-label settings, significantly expanding the label space.
This expansion limited to the capacity to provide all examples for each prompt. 
The lack of a standardized approach to these few-shot prompts would be a notable challenge.

\begin{table*}
    \small
    \centering
    \renewcommand{\arraystretch}{1.5}
    \setlength{\tabcolsep}{0.5em}
    
    \begin{tabular}{>{\raggedright\arraybackslash}p{15cm}}
        \toprule
        You are an Intent Detection Model on single utterance. \\
        
        \textbf{[Task Definition]} Detect single or more intent(s) of each utterance, but you can only classify UP TO 3 most plausible intents on 1 utterance.\\
        
        \textbf{[Intents]} \texttt{atis\_airport, atis\_ground\_service, atis\_abbreviation, atis\_city, atis\_aircraft, atis\_ground\_fare, atis\_flight, ...} \\
        
        \textbf{[Answer format]} If more than one, concatenate with `\#', such as \{Intent\}\#\{Intent\}. 
        
        e.g. \texttt{atis\_ground\_fare\#atis\_distance} \\
        \arrayrulecolor{lightgray}\midrule\arrayrulecolor{black} 
        
        % \textcolor{commentcolor}{\# $N$ iterations}\\
            \textbf{[Example 1]}
            \begin{quote}
                {[Utterance]} does delta aircraft fly dc10
                
                {[Answer]} \texttt{atis\_aircraft} 
            \end{quote}
            
            \textbf{[Example 2]}
            \begin{quote}
                {[Utterance]} which airline has more business class flights than any other airline and what city is the airport mco in
                
                {[Answer]} \texttt{atis\_airline\#atis\_city} 
            \end{quote}
            
            \textbf{[Example 3]}
            \begin{quote}
            {[Utterance]} what does the fare code qx mean, what is the distance between pittsburgh airport and downtown pittsburgh and what is restriction ap80
            
            {[Answer]} \texttt{atis\_abbreviation\#atis\_distance\#atis\_restriction} 
            \end{quote}
            \\
            \arrayrulecolor{lightgray}\midrule\arrayrulecolor{black}
        
        \textbf{[Query]} Detect a single or up to 3 intent(s) on this following utterance. : $\mathrm{utt}$ \\
        
        \bottomrule
    \end{tabular}
    
    \caption{Specification of the prompt used for addressing MID within the in-context learning framework. (Dataset: MixATIS).}
    \centering
    \label{table:appendix_chatgpt_prompt}
\end{table*}